\documentclass{article}
\usepackage[utf8]{inputenc}
\usepackage[spanish]{babel}
\usepackage{hyperref}
\usepackage{graphicx}
\usepackage{float}
\usepackage{booktabs}
\usepackage{geometry}
\usepackage{amsmath}
\geometry{margin=2.5cm}
\graphicspath{{imagenes/}}

\title{\textbf{Análisis de Acceso y Eficiencia del Sistema\\
Notarial Colombiano}\\
\large Un enfoque departamental con datos del SNR y el DANE}
\author{Juan Felipe Gutiérrez}
\date{Junio 2026}

\begin{document}

\maketitle
\tableofcontents
\newpage

% ─────────────────────────────────────────
\section{Introducción y razón del análisis}
% ─────────────────────────────────────────

Colombia cuenta con 32 departamentos y el Distrito Capital 
de Bogotá, cada uno con realidades económicas, demográficas 
y sociales profundamente distintas. Sin embargo, el sistema 
notarial opera con tarifas uniformes reguladas por la 
Superintendencia de Notariado y Registro (SNR), sin importar 
si el servicio se presta en Bogotá o en Vaupés.

Esta uniformidad regulatoria plantea una pregunta que hasta 
ahora no ha sido respondida con datos:

\begin{center}
\textit{¿El acceso a servicios notariales en Colombia es 
equitativo entre departamentos, o existen brechas 
sistemáticas donde la población tiene menor acceso 
relativo a estos servicios?}
\end{center}

Este proyecto construye un \textbf{índice de acceso y 
eficiencia notarial por departamento}, cruzando datos 
oficiales del SNR con variables demográficas y económicas 
del DANE y el DNP, para responder tres preguntas concretas:

\begin{enumerate}
    \item \textbf{¿Dónde hay más actos notariales de los 
    esperados} para el tamaño poblacional y económico del 
    departamento? — Departamentos sobre-servidos.
    \item \textbf{¿Dónde hay menos actos de los esperados?} 
    — Departamentos sub-servidos con posible brecha de acceso.
    \item \textbf{¿Qué variables explican mejor el volumen 
    notarial?} — Para entender qué mueve el mercado 
    notarial colombiano.
\end{enumerate}

\subsection{Relevancia e impacto del análisis}

\begin{itemize}
    \item \textbf{SNR y Estado colombiano:} Insumo para 
    política pública de distribución equitativa de notarías 
    y acceso a servicios formales.
    \item \textbf{Operadores del sector notarial:} Identifica 
    departamentos con demanda notarial insatisfecha.
    \item \textbf{Academia e investigación:} Primer análisis 
    cuantitativo del mercado notarial colombiano con datos 
    abiertos, replicable y actualizable anualmente.
    \item \textbf{Ciudadanos y sociedad civil:} Evidencia 
    sobre inequidades en el acceso a servicios jurídicos 
    formales entre regiones.
\end{itemize}

\newpage
% ─────────────────────────────────────────
\section{Marco de análisis de datos}
% ─────────────────────────────────────────

Este proyecto sigue el marco estándar de los cuatro tipos 
de análisis de datos, aplicados en secuencia:

\begin{enumerate}
    \item \textbf{Análisis descriptivo} 
    (\textit{¿qué sucedió?}): Volumen de actos notariales 
    por departamento, participación de mercado y distribución 
    geográfica. Provee la fotografía del estado actual del 
    sistema notarial colombiano.

    \item \textbf{Análisis diagnóstico} 
    (\textit{¿por qué sucedió?}): Construcción del índice 
    per cápita y cálculo de desviaciones respecto al promedio 
    nacional. Cruce con variables demográficas y económicas 
    del DANE para identificar causas de las brechas. 
    Regresiones y análisis de panel en Stata.

    \item \textbf{Análisis predictivo} 
    (\textit{¿qué podría pasar?}): Con la serie histórica 
    2021--2026 del SNR se proyecta la evolución del acceso 
    notarial por departamento. Modelos de series de tiempo 
    y regresión múltiple en Stata.

    \item \textbf{Análisis prescriptivo} 
    (\textit{¿qué debería hacerse?}): Recomendaciones de 
    política pública y estrategia basadas en los hallazgos, 
    visualizadas en Power BI.
\end{enumerate}

\subsection{Indicadores clave de desempeño (KPIs)}

\begin{table}[H]
\centering
\begin{tabular}{lll}
\toprule
\textbf{KPI} & \textbf{Fuente} & \textbf{Fase} \\
\midrule
Actos notariales per cápita & SNR + DANE & Descriptiva \\
Participación departamental (\%) & SNR & Descriptiva \\
Posición en ranking nacional & SNR & Descriptiva \\
Desviación vs.\ promedio nacional & SNR + DANE & Diagnóstica \\
Clasificación acceso notarial & Índice propio & Diagnóstica \\
Correlación PIB--volumen notarial & SNR + DANE & Diagnóstica \\
Correlación IPM--volumen notarial & SNR + DNP & Diagnóstica \\
Proyección de crecimiento (\%) & SNR histórico & Predictiva \\
Índice de brecha de acceso & Modelo Stata & Predictiva \\
\bottomrule
\end{tabular}
\caption{KPIs del análisis de acceso notarial colombiano}
\end{table}

\subsection{Metodología central: proporciones y desviaciones}

El índice central del análisis es el cociente de actos 
notariales per cápita por departamento:

\begin{equation}
\text{Actos per cápita}_i = 
\frac{\text{Total actos notariales}_i}{\text{Población}_i}
\end{equation}

La desviación respecto al promedio nacional se calcula como:

\begin{equation}
\text{Desviación}_i = 
\frac{\text{Actos per cápita}_i - \bar{x}}{\bar{x}} 
\times 100
\end{equation}

donde $\bar{x}$ es el promedio nacional ponderado por 
población. Un valor positivo indica un departamento 
\textbf{sobre-servido}; un valor negativo indica un 
departamento \textbf{sub-servido}.

\begin{table}[H]
\centering
\begin{tabular}{lll}
\toprule
\textbf{Desviación} & \textbf{Clasificación} & 
\textbf{Interpretación} \\
\midrule
$> +50\%$ & Sobre-servido & 
Muy por encima del promedio \\
$+10\%$ a $+50\%$ & Sobre promedio & 
Por encima del promedio \\
$-10\%$ a $+10\%$ & Promedio & 
En la media nacional \\
$-10\%$ a $-50\%$ & Sub-servido & 
Por debajo del promedio \\
$< -50\%$ & Brecha crítica & 
Muy por debajo del promedio \\
\bottomrule
\end{tabular}
\caption{Escala de clasificación de acceso notarial}
\end{table}

\newpage
% ─────────────────────────────────────────
\section{Fuentes de datos y pipeline}
% ─────────────────────────────────────────

El proyecto usa cuatro fuentes de datos públicas, 
procesadas automáticamente con Python:

\begin{figure}[H]
    \centering
    \includegraphics[width=\textwidth]{stack_tecnologico.png}
    \caption{Pipeline completo — fuentes, recolección, 
    procesamiento y análisis}
    \label{fig:pipeline}
\end{figure}

\subsection{Fuente 1 — SNR: Actos Jurídicos Notariales}

\textbf{Qué contiene:} Base de datos publicada en 
\texttt{datos.gov.co} con 40 tipos de actos notariales 
desagregados por departamento y fecha, con cobertura 
mensual desde enero de 2021 hasta abril de 2026.

\textbf{Cobertura:} 33 entidades territoriales 
(32 departamentos + Bogotá D.C.)\\
\textbf{Última actualización:} 5 de junio de 2026\\
\textbf{URL:} \url{https://www.datos.gov.co/Justicia-y-Derecho/Actos-Jur-dicos-Notariales/qktb-5f8m/about_data}

\subsection{Fuente 2 — DANE: Proyecciones de población}

\textbf{Qué contiene:} Proyecciones de población por 
departamento basadas en el Censo Nacional de Población 
y Vivienda 2018. Variable de normalización fundamental 
para calcular los actos per cápita.

\textbf{URL:} \url{https://www.dane.gov.co/index.php/estadisticas-por-tema/demografia-y-poblacion/proyecciones-de-poblacion}

\subsection{Fuente 3 — DANE: PIB departamental}

\textbf{Qué contiene:} Producto Interno Bruto por 
departamento a precios corrientes. Permite determinar 
si las brechas en acceso notarial se explican por 
diferencias en actividad económica.

\textbf{URL:} \url{https://www.dane.gov.co/index.php/estadisticas-por-tema/cuentas-nacionales/cuentas-nacionales-departamentales}

\subsection{Fuente 4 — DNP: Índice de Pobreza 
Multidimensional (IPM)}

\textbf{Qué contiene:} Medida de pobreza que combina 
múltiples dimensiones — educación, salud, vivienda, 
trabajo y condiciones de vida — por departamento. 
Permite evaluar si las brechas de acceso notarial 
se relacionan con condiciones de pobreza estructural.

\textbf{URL:} \url{https://www.datos.gov.co}\\
\textbf{Búsqueda:} \textit{``índice pobreza 
multidimensional departamento''}

\newpage
% ─────────────────────────────────────────
\section{Herramientas y stack tecnológico}
% ─────────────────────────────────────────

\subsection{Python con pandas}

\textbf{Para qué se usa:} Limpieza y transformación 
de los datos del SNR, DANE y DNP. Cálculo del índice 
per cápita, desviaciones y clasificaciones. Carga 
automática a SQLite.

\textbf{Librerías:}
\begin{itemize}
    \item \texttt{pandas} — manipulación de datos
    \item \texttt{matplotlib} / \texttt{seaborn} — 
    visualizaciones para el informe
    \item \texttt{sqlite3} — conexión a la base de datos
\end{itemize}

\subsection{SQLite y SQL}

\textbf{Para qué se usa:} Base de datos central del 
proyecto. Almacena todas las fuentes y permite 
consultas cruzadas entre ellas.

\textbf{Tablas principales:}
\begin{itemize}
    \item \texttt{actos\_notariales} — datos SNR
    \item \texttt{poblacion\_departamentos} — DANE
    \item \texttt{pib\_departamentos} — DANE
    \item \texttt{ipm\_departamentos} — DNP
    \item \texttt{indice\_acceso} — KPIs consolidados
\end{itemize}

\textbf{Consulta central del análisis:}
\begin{verbatim}
SELECT a.Departamento,
       a.total_actos,
       p.poblacion,
       ROUND(a.total_actos * 1.0 
           / p.poblacion, 4) AS actos_per_capita,
       ROUND((a.total_actos * 1.0 / p.poblacion 
           - avg_nacional) / avg_nacional 
           * 100, 2) AS desviacion_pct
FROM resumen_departamento a
JOIN poblacion_departamentos p 
    ON a.Departamento = p.Departamento
ORDER BY desviacion_pct DESC;
\end{verbatim}

\subsection{Stata}

\textbf{Modelo de regresión principal:}

\begin{equation}
\text{actos\_per\_cápita}_i = \beta_0 + 
\beta_1 \cdot \text{PIB\_per\_cápita}_i + 
\beta_2 \cdot \text{urbanización}_i + 
\beta_3 \cdot \text{IPM}_i + 
\beta_4 \cdot \text{num\_notarías}_i + 
\varepsilon_i
\end{equation}

\textbf{Análisis:}
\begin{itemize}
    \item Regresión múltiple por corte transversal
    \item Análisis de panel 2021--2026
    \item Series de tiempo por departamento
    \item Tests de significancia estadística
\end{itemize}

\subsection{Power BI}

Dashboard interactivo con mapa de Colombia coloreado 
por índice de acceso notarial, ranking departamental 
y evolución temporal 2021--2026. Conectado directamente 
a la base de datos SQLite.

\subsection{LaTeX con Tectonic}

Informe final compilado desde VS Code con actualización 
automática. Versionado en GitHub.

\textbf{Repositorio:}
\url{https://github.com/juanguti1231/notaria74-benchmarking-comparacion-con-competencia-}

\newpage
% ─────────────────────────────────────────
\section{Resultados descriptivos preliminares}
% ─────────────────────────────────────────

\subsection{Cobertura del dataset}

El dataset del SNR cubre enero 2021 -- abril 2026, 
con datos mensuales para las 33 entidades territoriales 
de Colombia. El total supera las 57,000 filas con 
40 variables de tipos de actos notariales.

\subsection{Participación de Bogotá en compraventas}

\begin{table}[H]
\centering
\begin{tabular}{cc}
\toprule
\textbf{Año} & \textbf{Participación (\%)} \\
\midrule
2021 & 15.22 \\
2022 & 15.95 \\
2023 & 14.87 \\
2024 & 15.08 \\
2025 & 15.00 \\
2026 & 16.06 \\
\bottomrule
\end{tabular}
\caption{Participación de Bogotá D.C. en compraventas 
notariales nacionales 2021--2026}
\end{table}

\subsection{Ranking departamental 2025}

\begin{table}[H]
\centering
\begin{tabular}{clrr}
\toprule
\textbf{Pos.} & \textbf{Departamento} & 
\textbf{Total actos} & \textbf{Compraventas} \\
\midrule
1  & Bogotá D.C.     & 12,492,643 &  96,795 \\
2  & Antioquia       &  6,250,292 & 112,231 \\
3  & Valle del Cauca &  3,958,684 &  54,043 \\
4  & Cundinamarca    &  1,897,069 &  40,906 \\
5  & Santander       &  1,815,181 &  40,719 \\
6  & Atlántico       &  1,741,451 &  18,510 \\
7  & Bolívar         &  1,284,247 &  13,152 \\
8  & Meta            &  1,274,090 &  23,855 \\
9  & Risaralda       &  1,181,143 &  17,435 \\
10 & Tolima          &  1,044,383 &  29,362 \\
\bottomrule
\end{tabular}
\caption{Ranking departamental por total de actos 
notariales 2025}
\end{table}

\subsection{Hallazgo clave: Bogotá vs Antioquia}

Aunque Bogotá lidera en total de actos (12.49M vs 6.25M), 
\textbf{Antioquia supera a Bogotá en compraventas} 
(112,231 vs 96,795), lo que sugiere que el mercado 
inmobiliario antioqueño es proporcionalmente más activo. 
Bogotá concentra su actividad principalmente en 
\textbf{autenticaciones} (10.89 millones — 87\% 
de su total).

\subsection{Top de actos en Bogotá 2025}

\begin{table}[H]
\centering
\begin{tabular}{lr}
\toprule
\textbf{Tipo de acto} & \textbf{Cantidad} \\
\midrule
Autenticaciones              & 10,897,340 \\
Copias del Registro Civil    &    882,159 \\
Declaraciones Extra Juicio   &    226,332 \\
Compraventas                 &     96,795 \\
Cancelaciones de Hipotecas   &     73,301 \\
Registros Civiles Nacimiento &     66,615 \\
Registros Civiles Defunción  &     60,162 \\
Otras --- Bien Inmueble      &     40,193 \\
Vivienda de Interés Social   &     38,748 \\
Registros Civiles Matrimonio &     23,150 \\
\bottomrule
\end{tabular}
\caption{Top 10 actos notariales en Bogotá D.C. 2025}
\end{table}

\subsection{Próximos pasos}

\begin{enumerate}
    \item Descargar población y PIB por departamento 
    del DANE e IPM del DNP
    \item Calcular índice de actos per cápita 
    por departamento
    \item Calcular desviaciones respecto al 
    promedio nacional
    \item Clasificar departamentos según escala 
    de acceso notarial
    \item Correr modelo de regresión en Stata
    \item Construir dashboard en Power BI
\end{enumerate}

\end{document}